\documentclass[twoside,final]{hcmut-report}

\usepackage{listings}
\usepackage{xcolor}
\usepackage{tcolorbox}
\usepackage[final]{microtype}
\usepackage{xurl}

\setlength{\emergencystretch}{3em}
\Urlmuskip=0mu plus 1mu

\lstdefinelanguage{PowerShell}{
    keywords={param,if,else,elseif,foreach,for,while,function,return,throw,Write-Host,Set-Content},
    sensitive=false,
    morecomment=[l]{\#},
    morestring=[b]",
    morestring=[b]'
}

% ─── Python code listing style ────────────────────────────────────────────
\definecolor{codebg}{rgb}{0.96,0.96,0.96}
\definecolor{codekey}{rgb}{0.13,0.29,0.53}
\definecolor{codestr}{rgb}{0.58,0.00,0.10}
\definecolor{codecmt}{rgb}{0.40,0.40,0.40}
\lstdefinestyle{pystyle}{
    backgroundcolor=\color{codebg},
    basicstyle=\ttfamily\footnotesize,
    keywordstyle=\color{codekey}\bfseries,
    stringstyle=\color{codestr},
    commentstyle=\color{codecmt}\itshape,
    language=Python,
    frame=single,
    breaklines=true,
    showstringspaces=false,
    tabsize=2,
    captionpos=b,
    numbers=left,
    numberstyle=\tiny\color{codecmt}
}
\lstset{style=pystyle}

\coursename{Software Testing}
\reporttype{Project 3:}
\title{Data-driven Automation Testing and Non-functional Testing}
\advisor{& Bùi Hoài Thắng &}
\groupname{& 11 &}
\stuname{%
  & Lâm Minh Tùng Quân &  2352991 \\
  & Trương Xuân Sang &   2353045\\
  & Trương Gia Kỳ Nam &  2352787 \\
  & Nguyễn Hữu Phúc &  2352938 \\
  & Châu Anh Nhật  & 2352856\\
}

\begin{document}

\coverpage%

\tableofcontents
\pagebreak

% ============================================================
% CONTRIBUTION TABLE
% ============================================================
\begin{table}[ht]
    \centering
    \caption{Project Contribution and Workload Distribution}
    \vspace{0.2cm}
    \small
    \renewcommand{\arraystretch}{1.4}

    \begin{tabularx}{\textwidth}{|c|c|p{2.5cm}|>{\raggedright\arraybackslash}X|c|}
    \hline
    \textbf{No.} & \textbf{ID} & \textbf{Full Name} & \textbf{Description} & \textbf{\%} \\ \hline

    1 & 2352938 & Nguyễn Hữu Phúc &
    \begin{filelist}
        \item Feature 001: Admin Adds a New User -- Level 1 + Level 2 automation
        \item NFR: Performance Testing -- Locust load test + SLA checks (\texttt{TC-001.py})
    \end{filelist} & 20\% \\ \hline

    2 & 2352787 & Trương Gia Kỳ Nam &
    \begin{filelist}
       \item Feature 006: Teacher Sets Up a Quiz -- Level 1 + Level 2 automation
       \item NFR: Usability Testing -- keyboard navigation (\texttt{TC-006.py})
    \end{filelist} & 20\% \\ \hline

    3 & 2352856 & Châu Anh Nhật &
    \begin{filelist}
        \item Feature 005: User Creates Calendar Event -- Level 1 + Level 2 automation
        \item NFR: Compatibility Testing -- responsive viewport (\texttt{TC-005.py})
    \end{filelist} & 20\% \\ \hline

    4 & 2353045 & Trương Xuân Sang &
    \begin{filelist}
       \item Feature 002: Admin Creates a New Course -- Level 1 + Level 2 automation
       \item NFR: Security Testing -- passive HTTP probes (\texttt{TC-002.py})
    \end{filelist} & 20\% \\ \hline

    5 & 2352991 & Lâm Minh Tùng Quân &
    \begin{filelist}
        \item Feature 003: Teacher Creates Assignment -- Level 1 + Level 2 automation
       \item Feature 004: Teacher Grades an Assignment -- Level 1 + Level 2 automation
       \item NFR: Accessibility Testing -- axe-core WCAG 2.1 audit (\texttt{TC-003.py})
        \item NFR: Reliability Testing -- repeated-request consistency (\texttt{TC-004.py})
    \end{filelist} & 20\% \\ \hline

    \end{tabularx}
\end{table}
\newpage

% ============================================================
% 1. GENERAL INFORMATION
% ============================================================
\section{General Information}

\subsection{Project Overview}
Project~\#3 extends the manual black-box test designs from Project~\#2 by transforming them into \textbf{data-driven automated test suites} written in Python with the Selenium WebDriver framework. In addition to the functional regression suites, this project introduces \textbf{non-functional testing} covering six quality attributes -- Performance, Security, Accessibility, Reliability, Compatibility, and Usability -- of the Moodle LMS demo environment.

The complete artefact set consists of:
\begin{itemize}
    \item \textbf{Level 1 automation} (data only from CSV) -- six suites, 169 parametric test methods.
    \item \textbf{Level 2 automation} (data \emph{and} locators / URLs from CSV) -- a unified suite, 169 parametric test methods.
    \item \textbf{Non-functional testing} -- six Python files each implementing a distinct NFR technique (Performance, Security, Accessibility, Reliability, Compatibility, Usability), one file scoped to each of the six functional features.
    \item \textbf{Supporting tooling} -- a master PowerShell runner (\texttt{run\_all.ps1}), a bash runner (\texttt{run\_all.sh}), a Selenium-based test-data cleanup script (\texttt{cleanup\_moodle.py}), and a pinned \texttt{requirements.txt}.
\end{itemize}

\subsection{Target Application}
All automation runs against the dedicated MoodleCloud LMS instance hosted at
\begin{center}
    \textbf{\url{https://xuansang1234.moodlecloud.com/}}
\end{center}
which was established during Project~\#2 to bypass the instability of the public Moodle demo. The application instance, the administrator credentials, and the six feature flows under test remain identical to those of Project~\#2; only the testing technique has evolved.

\subsection{Test Environment}
To ensure reproducibility, every test session was conducted under the following pinned configuration:

\begin{table}[H]
\centering
\renewcommand{\arraystretch}{1.3}
\begin{tabularx}{\textwidth}{|l|X|}
\hline
\textbf{Component} & \textbf{Version / Detail} \\ \hline
Operating System & Windows 11 Pro 64-bit (build 22631) \\ \hline
Language & Python 3.10.11 (CPython) \\ \hline
Browser & Google Chrome 148.x (auto-managed ChromeDriver) \\ \hline
Test framework & Selenium 4.10+, pytest 9.0, unittest \\ \hline
Driver management & \texttt{webdriver-manager} 4.0 (auto ChromeDriver sync) \\ \hline
Data libraries & \texttt{pandas} 2.0, \texttt{openpyxl} 3.1 \\ \hline
Load testing & \texttt{locust} 2.43 \\ \hline
HTTP testing & \texttt{requests} 2.31 (passive security probes \& reliability checks) \\ \hline
Accessibility audit & \texttt{axe-selenium-python} 2.1 (axe-core engine) \\ \hline
Display resolution & 1920~$\times$~1080 (Full HD) at 100\% scale \\ \hline
\end{tabularx}
\caption{Pinned test environment}
\end{table}

% ============================================================
% 2. REPOSITORY LAYOUT
% ============================================================
\section{Repository Layout and Submission Package}

The submission archive is organised into three top-level folders matching the project rubric. All test files follow the unified \texttt{TC-XXX*} naming convention so the relationship between feature, code, and data is immediate at a glance.

\begin{lstlisting}[language=bash, basicstyle=\ttfamily\footnotesize]
swe-testing/
|-- level1/                          # Level 1 - data only from CSV
|   |-- TC-001_code.py               # Admin Adds User       (43 rows)
|   |-- TC-001_data.csv
|   |-- TC-002_code.py               # Admin Creates Course  (27 rows)
|   |-- TC-002_data.csv
|   |-- TC-003_code.py               # Teacher Creates Assignment (27 rows)
|   |-- TC-003_data.csv
|   |-- TC-004_code.py               # Teacher Grades Student (17 rows)
|   |-- TC-004_data.csv
|   |-- TC-005_code.py               # Admin Creates Calendar Event (28 rows)
|   |-- TC-005_data.csv
|   |-- TC-006_code.py               # Teacher Sets Up a Quiz (27 rows)
|   `-- TC-006_data.csv
|
|-- level2/                          # Level 2 - data + locators from CSV
|   |-- TC-ALL.py                    # Single unified script (169 tests)
|   |-- TC-001_data.csv
|   |-- TC-002_data.csv
|   |-- TC-003_data.csv
|   |-- TC-004_data.csv
|   |-- TC-005_data.csv
|   `-- TC-006_data.csv
|
|-- non_functional/                  # Non-functional test suites
|   |-- TC-001.py                    # Add User       -> Performance (Locust)
|   |-- TC-002.py                    # Create Course  -> Security (requests)
|   |-- TC-003.py                    # Create Assign  -> Accessibility (axe-core)
|   |-- TC-004.py                    # Grade Student  -> Reliability (requests)
|   |-- TC-005.py                    # Calendar Event -> Compatibility (Selenium)
|   `-- TC-006.py                    # Quiz Setup     -> Usability (Selenium)
|
|-- pytest.ini                       # Discovery rules for TC-* file names
|-- requirements.txt                 # Pinned dependencies
|-- run_all.ps1                      # Master PowerShell runner (Windows)
|-- run_all.sh                       # Master bash runner (macOS / Linux)
|-- cleanup_moodle.py                # Test-data cleanup utility
|-- README.md                        # Full usage guide
\end{lstlisting}

% ============================================================
% 3. LEVEL 1 -- DATA-DRIVEN AUTOMATION
% ============================================================
\section{Level 1 -- Data-driven Automation Testing}

\subsection{Design Principles}
Level~1 implements the data-driven testing pattern as defined in the project description: each test case is parametrised by a row of input data and an \texttt{expected\_result} value read from a CSV file, while the locators (element identifiers) remain hard-coded inside the Python script.

The architecture for every Level~1 test file follows the same template:
\begin{enumerate}
    \item A module-level CSV reader (\texttt{csv.DictReader}) loads all rows from \texttt{TC-XXX\_data.csv}.
    \item A factory function (\texttt{\_make\_test(row)}) returns a closure that becomes the body of one \texttt{unittest} test method.
    \item A \texttt{setattr} loop attaches one such method per CSV row at import time, so \texttt{pytest} collects them as independent tests.
    \item Each test method calls the same \texttt{\_fill\_and\_submit(row)} helper, then asserts that the observed outcome matches \texttt{row["expected\_result"]}.
\end{enumerate}

A repository-level \texttt{pytest.ini} sets \texttt{python\_files = TC-*.py test\_*.py} and \texttt{addopts = --import-mode=importlib} so the hyphenated file names (\texttt{TC-001\_code.py}, etc.) are discovered automatically.

\subsection{Per-feature Implementation Details}

This section documents each of the six features automated in Level~1. Test
case purposes are grouped by the black-box technique that generated them
(BVA, ECP, DT, UC) -- mirroring the categorisation used in the Project~\#2
report. Each CSV row in the data file corresponds to one parametric test
method named \texttt{test\_TC\_XXX\_YYY} at pytest collection time.

\begin{tcolorbox}[colback=yellow!8!white, colframe=orange!70!black, title=Note on \textbf{Exp.} column values, fonttitle=\small\bfseries, left=4pt, right=4pt, top=4pt, bottom=4pt]
\small
In the tables below, the \textbf{Exp.}\ column uses shorthand labels (\texttt{success}, \texttt{fail}, \texttt{success/fail}, \texttt{varies}).
The actual \texttt{expected\_result} values stored in each CSV file are the \emph{exact Moodle error message strings} returned by the application -- for example \texttt{"Passwords must be at least 8 characters long"}, \texttt{"Required"}, \texttt{"This username already exists"}, \texttt{"The course end date must be after the start date"}, etc.
The test script matches these strings as case-insensitive substrings of the page source; \texttt{success} rows contain the word \texttt{"success"} and trigger a positive assertion, while all other rows assert that the corresponding error message appears in the rendered HTML.
\end{tcolorbox}

% ─────────────────────────────────────────────────────────────────────────
\subsubsection{Feature 001 -- Admin Adds a New User (43 test cases)}

\textbf{Target URL:} \url{https://xuansang1234.moodlecloud.com/user/editadvanced.php?id=-1}\\
\textbf{Pre-condition:} Logged in as administrator on the Moodle site.\\
\textbf{Data file:} \texttt{level1/TC-001\_data.csv} (43 rows).\\
\textbf{CSV schema:} \texttt{test\_case\_id, username, password, firstname, lastname, email, expected\_result}.

\textbf{Special handling:}
\begin{itemize}
    \item Moodle marks \texttt{\#id\_newpassword} as \texttt{readonly}; the script bypasses this through JavaScript injection (\texttt{removeAttribute}~+~\texttt{dispatchEvent}).
    \item The sentinel value \texttt{\_\_generate\_\_} in the password column ticks the \emph{Generate password and notify user} checkbox instead of typing a password.
\end{itemize}
\textbf{Outcome detection:} success $\Leftrightarrow$ ``Changes saved'' substring in page source; fail $\Leftrightarrow$ any \texttt{[id\^{}='id\_error\_']} element rendered.

\textbf{Test-case purpose grouped by technique:}
\begin{table}[H]
\centering
\renewcommand{\arraystretch}{1.2}
\begin{tabularx}{\textwidth}{|c|l|X|c|}
\hline
\textbf{TC range} & \textbf{Method} & \textbf{Boundary / scenario tested} & \textbf{Exp.} \\ \hline
TC-001-001 & BVA & Username (nom 17 chars) -- valid user, all fields nominal & success \\ \hline
TC-001-002 & BVA & Username (min$-$) empty string -- missing required field & fail \\ \hline
TC-001-003..004 & BVA & Username (min 1 char, min$+$ 2 chars) -- boundary lower & success \\ \hline
TC-001-005 & BVA & Password (min$-$ 7 chars) -- below 8-char minimum & fail \\ \hline
TC-001-006..010 & BVA & Password length boundaries (Moodle-documented min 8; min$+$1=9, max$-$=127, max=128, max$+$=129 chars) & success/fail \\ \hline
TC-001-011 & BVA & First name (min$-$) empty -- required field missing & fail \\ \hline
TC-001-012..015 & BVA & First name boundaries (1, 2, 99, 100 chars) & success \\ \hline
TC-001-016 & BVA & Last name (min$-$) empty -- required field missing & fail \\ \hline
TC-001-017..020 & BVA & Last name boundaries (1, 2, 99, 100 chars) & success \\ \hline
TC-001-021..024 & ECP & Invalid username formats: uppercase, space, special char, duplicate of TC-001-001 username & fail \\ \hline
TC-001-025..028 & DT & Password complexity rules: missing digit / missing upper / missing lower / numeric-sequence detected & fail \\ \hline
TC-001-029..031 & ECP & Invalid email formats: no @, missing domain, embedded space & fail \\ \hline
TC-001-032 & UC & Happy-path: all fields nominal -- baseline success scenario & success \\ \hline
TC-001-033 & UC & \texttt{\_\_generate\_\_} sentinel -- alternate flow with auto-password & success \\ \hline
TC-001-034 & UC & Duplicate username (same as TC-001-001 and TC-001-024) -- second attempt & fail \\ \hline
TC-001-035..043 & ECP/DT & Additional password \& email edge combinations & varies \\ \hline
\end{tabularx}
\caption{TC-001 test case purpose by group (43 rows total)}
\end{table}

% ─────────────────────────────────────────────────────────────────────────
\subsubsection{Feature 002 -- Admin Creates a New Course (27 test cases)}

\textbf{Target URL:} \url{https://xuansang1234.moodlecloud.com/course/edit.php?category=0}\\
\textbf{Data file:} \texttt{level1/TC-002\_data.csv} (27 rows).\\
\textbf{CSV schema:} \texttt{test\_case\_id, fullname, shortname, end\_date\_enabled, end\_date\_offset\_days, end\_date\_offset\_years, numsections, expected\_result}.

\textbf{Special handling:} End-date day / month / year drop-downs are set through JavaScript helpers (\texttt{sS()} for select boxes, \texttt{ens()}~/~\texttt{dis()} for the \texttt{id\_enddate\_enabled} checkbox).\\
\textbf{Outcome detection:} success $\Leftrightarrow$ ``Announcements'' substring (course landing page); fail $\Leftrightarrow$ any \texttt{[id\^{}='id\_error\_']} validation element.

\textbf{Test-case purpose grouped by technique:}
\begin{table}[H]
\centering
\renewcommand{\arraystretch}{1.2}
\begin{tabularx}{\textwidth}{|c|l|X|c|}
\hline
\textbf{TC range} & \textbf{Method} & \textbf{Boundary / scenario tested} & \textbf{Exp.} \\ \hline
TC-002-001 & BVA & Full name (nom 20 chars) + shortname \texttt{sn001\_xxxxxxxx} -- nominal success & success \\ \hline
TC-002-002 & BVA & Full name (min$-$) empty -- required field missing & fail \\ \hline
TC-002-003..004 & BVA & Full name boundaries (1, 2 chars) & success \\ \hline
TC-002-005..006 & BVA & Full name upper boundaries (max$-$=254, max=255 chars; UI enforces limit, max$+$1 not applicable) & success \\ \hline
TC-002-007 & BVA & Short name (min$-$) empty -- required field missing & fail \\ \hline
TC-002-008..009 & BVA & Short name boundaries (1, 2 chars) & success \\ \hline
TC-002-010..011 & BVA & Short name upper boundaries (max$-$=254, max=255 chars; UI enforces limit, max$+$1 not applicable) & success \\ \hline
TC-002-012..013 & BVA & End date chronology (-1 day, today vs start) -- invalid order & fail \\ \hline
TC-002-014..017 & BVA & End date boundaries (+1 day .. +12 years from start) & success \\ \hline
TC-002-018 & ECP & Duplicate short name -- second create attempt with same shortname & fail \\ \hline
TC-002-019..023 & DT & Course format \texttt{weekly}/\texttt{topics} + end-date enabled/disabled combinations & success \\ \hline
TC-002-024..025 & UC & Happy paths (all valid + end-date disabled variants) & success \\ \hline
TC-002-026 & UC & Duplicate short name as intentional second attempt -- mirrors TC-002-018 & fail \\ \hline
TC-002-027 & UC & Past end date -- date-conflict exception scenario & fail \\ \hline
\end{tabularx}
\caption{TC-002 test case purpose by group (27 rows total)}
\end{table}

% ─────────────────────────────────────────────────────────────────────────
\subsubsection{Feature 003 -- Teacher Creates Assignment (27 test cases)}

\textbf{Target URL:} \url{https://xuansang1234.moodlecloud.com/course/modedit.php?add=assign&type&course=10&sectionid=39&return=0&beforemod=0}\\
\textbf{Role switch:} \texttt{setUpClass()} navigates to \path|/course/switchrole.php?id=1&switchrole=-1&returnurl=%2Fmy%2Findex.php| and clicks the Teacher button to switch the active viewing role to Teacher before any assignment tests run.\\
\textbf{Data file:} \texttt{level1/TC-003\_data.csv} (27 rows).\\
\textbf{CSV schema:} \path|test_case_id, name, gradepass, duedate_enabled, duedate_offset_days, duedate_offset_years, cutoff_offset_days, cutoff_offset_years, submission_file, submission_onlinetext, expected_result|.

\textbf{Special handling:} Three date fields (\texttt{allowsubmissionsfromdate}, \texttt{duedate}, \texttt{cutoffdate}) are populated via a single \texttt{execute\_script()} JS helper that also enables / disables their checkboxes per CSV row. Submission-type checkboxes (\texttt{file}, \texttt{onlinetext}) are only changed when the CSV column value is not \texttt{default}.\\
\textbf{Outcome detection:} ``Announcements'' substring after redirect $\Rightarrow$ success; any \texttt{[id\^{}='id\_error\_']} $\Rightarrow$ fail.

\textbf{Test-case purpose grouped by technique:}
\begin{table}[H]
\centering
\renewcommand{\arraystretch}{1.2}
\begin{tabularx}{\textwidth}{|c|l|X|c|}
\hline
\textbf{TC range} & \textbf{Method} & \textbf{Boundary / scenario tested} & \textbf{Exp.} \\ \hline
TC-003-001 & BVA & Assignment name (nom 29 chars) -- baseline valid create & success \\ \hline
TC-003-002 & BVA & Name (min$-$) empty -- required field missing & fail \\ \hline
TC-003-003..004 & BVA & Name boundaries (1, 2 chars lower bound) & success \\ \hline
TC-003-005..006 & BVA & Name upper boundaries (255, 256 chars; UI-enforced max) & success \\ \hline
TC-003-007..009 & BVA & Grade-to-pass boundaries (99, 100, 101) -- max grade is 100 & success/fail \\ \hline
TC-003-010..014 & BVA & Due date relative to allow-submission date (today fail, +1, +2 days, +10/11 years) & success/fail \\ \hline
TC-003-015..019 & BVA & Cut-off date relative to due date ($-1$ day fail, same day, $+1$ day, $+10$/$11$ years) & success/fail \\ \hline
TC-003-020..023 & DT & Submission-type combinations: both, file only, online only, neither & success \\ \hline
TC-003-024..025 & UC & Happy paths (nominal, due-date disabled) & success \\ \hline
TC-003-026..027 & UC & Exception flows: empty name, past due date & fail \\ \hline
\end{tabularx}
\caption{TC-003 test case purpose by group (27 rows total)}
\end{table}

% ─────────────────────────────────────────────────────────────────────────
\subsubsection{Feature 004 -- Teacher Grades an Assignment (17 test cases)}

\textbf{Target URL:} \url{https://xuansang1234.moodlecloud.com/mod/assign/view.php?id=41&action=grader&userid=2}\\
\textbf{Pre-condition:} Logged in as administrator (with full grading rights on course~10).\\
\textbf{Data file:} \texttt{level1/TC-004\_data.csv} (17 rows).\\
\textbf{CSV schema:} \texttt{test\_case\_id, grade, expected\_result}.

\textbf{Special handling:} The grade input is a React-controlled component, so a regular \texttt{send\_keys} does not register. The script uses the native \path|HTMLInputElement.prototype.value| setter and dispatches \texttt{input}~/~\texttt{change}~/~\texttt{blur} events to make React re-render. After submit, a sentinel DOM element \texttt{\#\_\_test\_marker[data-has-error]} is injected via JavaScript so the outcome can be read after the page navigates away.

\textbf{Test-case purpose grouped by technique:}
\begin{table}[H]
\centering
\renewcommand{\arraystretch}{1.2}
\begin{tabularx}{\textwidth}{|c|l|X|c|}
\hline
\textbf{TC range} & \textbf{Method} & \textbf{Boundary / scenario tested} & \textbf{Exp.} \\ \hline
TC-004-001 & UC & Grade nominal (50) -- baseline success scenario & success \\ \hline
TC-004-002..004 & BVA & Grade lower-bound boundaries ($-0.01$ fail, 0 min, 0.01 min$+$) & fail/success \\ \hline
TC-004-005..006 & BVA & Grade upper-bound boundaries (99.99 max$-$, 100 max) & success \\ \hline
TC-004-007..009 & BVA & Grade out of range (100.01, $-5$, 105) -- above max / below min & fail \\ \hline
TC-004-010..012 & ECP & Non-numeric / empty grade (letters, special chars, empty) & fail/success \\ \hline
TC-004-013..017 & UC & Use-case scenario variants: repeat happy-path grade (50) $\times$2 (rows 013--014), empty grade, above-max (105), non-numeric & varies \\ \hline
\end{tabularx}
\caption{TC-004 test case purpose by group (17 rows total)}
\end{table}

% ─────────────────────────────────────────────────────────────────────────
\subsubsection{Feature 005 -- User Creates Calendar Event (28 test cases)}

\textbf{Target URL:} \url{https://xuansang1234.moodlecloud.com/calendar/view.php?view=month}\\
\textbf{Data file:} \texttt{level1/TC-005\_data.csv} (28 rows).\\
\textbf{CSV schema:} \path|test_case_id, name, duration_type, minutes, until_offset_days, repeat, expected_result|.

\textbf{Special handling:} Three duration modes (\texttt{none}, \texttt{minutes}, \texttt{until}) map to radio buttons (values 0, 1, 2). \texttt{minutes} additionally fills a numeric field; \texttt{until} sets a future date through the \texttt{sS()} helper offset by \texttt{until\_offset\_days}. The modal submit button is located via \texttt{//div[@role='dialog']//button[@data-action='save']}.\\
\textbf{Outcome detection:} modal dialog closes within 8~s + ``Calendar'' substring in page source.

\textbf{Test-case purpose grouped by technique:}
\begin{table}[H]
\centering
\renewcommand{\arraystretch}{1.2}
\begin{tabularx}{\textwidth}{|c|l|X|c|}
\hline
\textbf{TC range} & \textbf{Method} & \textbf{Boundary / scenario tested} & \textbf{Exp.} \\ \hline
TC-005-001..004 & BVA & Event title boundaries (nom, empty, 1, 2 chars) & success/fail \\ \hline
TC-005-005..010 & BVA & Duration in minutes boundaries (0, 1, 2, 9999998, 9999999, 10000000) & success/fail \\ \hline
TC-005-011..015 & BVA & ``Until'' date offset ($-1$, 0, $+1$, $+364$, $+365$ days) & success/fail \\ \hline
TC-005-016..018 & ECP & Invalid minutes inputs (1.5 decimal, alphabetic, negative) & fail \\ \hline
TC-005-019..021 & DT & Duration mode combinations (none, until only, minutes only) & success \\ \hline
TC-005-022..024 & DT & Repeat checkbox enabled in each duration mode & success \\ \hline
TC-005-025..028 & UC & Happy path + exception flows (empty title, negative duration) & varies \\ \hline
\end{tabularx}
\caption{TC-005 test case purpose by group (28 rows total)}
\end{table}

% ─────────────────────────────────────────────────────────────────────────
\subsubsection{Feature 006 -- Teacher Sets Up a Quiz (27 test cases)}

\textbf{Target URL:} \url{https://xuansang1234.moodlecloud.com/course/modedit.php?add=quiz&type&course=12&sectionid=39&return=0&beforemod=0}\\
\textbf{Role switch:} \texttt{setUpClass()} switches to Teacher role on course~12, with editing mode enabled.\\
\textbf{Data file:} \texttt{level1/TC-006\_data.csv} (27 rows).\\
\textbf{CSV schema:} \texttt{test\_case\_id, name, timeclose\_enabled, close\_offset\_days, close\_offset\_years, timelimit\_enabled, timelimit\_number, gradepass, expected\_result}.

\textbf{Special handling:} The \texttt{expected\_result} column contains the full Moodle error message string returned after form submission. The script matches this as a case-insensitive substring in the page source, giving precise per-error assertions.

\textbf{Test-case purpose grouped by technique:}
\begin{table}[H]
\centering
\renewcommand{\arraystretch}{1.2}
\begin{tabularx}{\textwidth}{|c|l|X|c|}
\hline
\textbf{TC range} & \textbf{Method} & \textbf{Boundary / scenario tested} & \textbf{Exp.} \\ \hline
TC-006-001..004 & BVA & Grade-to-pass boundaries (5, 9.99, 10, 10.01) -- above max grade & success/fail\_grade \\ \hline
TC-006-005..010 & BVA & Time-limit boundaries ($-1$, 0, 1, 998, 999, 1000) & success/fail\_time \\ \hline
TC-006-011..016 & BVA & Close-date boundaries ($-1$ day, today, $+1$ day, $+10$/$11$/$12$ years) & success/fail\_date \\ \hline
TC-006-017 & ECP & Empty quiz name -- required field missing & fail\_name \\ \hline
TC-006-018..024 & DT & Timing combinations (close-date and time-limit enable/disable, 7 rule variants) & success \\ \hline
TC-006-025..026 & UC & Composite exception flows (past close date $+$ grade $>$ max) & fail \\ \hline
TC-006-027 & UC & Empty name in happy-path context -- final exception scenario & fail\_name \\ \hline
\end{tabularx}
\caption{TC-006 test case purpose by group (27 rows total)}
\end{table}

\subsection{Sample Code Skeleton}
\begin{lstlisting}[caption={Generic Level~1 test-factory pattern},label={lst:lvl1-factory}]
import csv, os, unittest
from selenium import webdriver
from selenium.webdriver.common.by import By

CSV_PATH = os.path.join(os.path.dirname(__file__), "TC-001_data.csv")

def _load_rows():
    with open(CSV_PATH, newline="", encoding="utf-8") as f:
        return list(csv.DictReader(f))

class TestAddUserLevel1(unittest.TestCase):
    @classmethod
    def setUpClass(cls):
        cls.driver = webdriver.Chrome()
        cls._login()

    def _fill_and_submit(self, row):
        # Locators are hardcoded here (Level 1 contract)
        self.driver.find_element(By.ID, "username").send_keys(row["username"])
        # ... fill other fields ...
        self.driver.find_element(By.ID, "id_submitbutton").click()

    def _get_outcome(self):
        if "Changes saved" in self.driver.page_source:
            return "success"
        if self.driver.find_elements(By.CSS_SELECTOR, "[id^='id_error_']"):
            return "fail"
        return "unknown"

def _make_test(row):
    def test_method(self):
        self._fill_and_submit(row)
        self.assertEqual(self._get_outcome(),
                         row["expected_result"].strip().lower())
    test_method.__name__ = f"test_{row['test_case_id'].replace('-', '_')}"
    return test_method

for _r in _load_rows():
    setattr(TestAddUserLevel1, f"test_{_r['test_case_id'].replace('-','_')}",
            _make_test(_r))
\end{lstlisting}

% ============================================================
% 4. LEVEL 2 -- LOCATORS + URLS FROM CSV
% ============================================================
\section{Level 2 -- Locators and URLs Externalised}

\subsection{Design Principles}
Level~2 fully realises the project description's strongest requirement: \emph{``testing data and testing items such as site URLs and elements like text fields and buttons are provided to the script.''} In other words, the Python source contains \textbf{zero site-specific values}: every URL, every locator strategy, and every locator value is read from CSV columns.

\subsection{Locator Resolution Helper}
A single helper \texttt{loc(row, prefix)} translates a pair of CSV columns into a Selenium locator tuple:

\begin{lstlisting}[caption={Locator resolution from CSV},label={lst:loc-helper}]
from selenium.webdriver.common.by import By

BY_MAP = {
    "id": By.ID, "name": By.NAME, "css selector": By.CSS_SELECTOR,
    "xpath": By.XPATH, "link text": By.LINK_TEXT,
    "partial link text": By.PARTIAL_LINK_TEXT,
    "tag name": By.TAG_NAME, "class name": By.CLASS_NAME,
}

def loc(row, prefix):
    """Resolve (By.X, "value") from columns <prefix>_locator_type and
       <prefix>_locator_value."""
    return (BY_MAP[row[f"{prefix}_locator_type"].strip().lower()],
            row[f"{prefix}_locator_value"].strip())
\end{lstlisting}

Every interactive element thus contributes two CSV columns, e.g.\ for the \texttt{username} input on TC-001 Level~2:
\begin{center}
    \texttt{username\_locator\_type = id, \quad username\_locator\_value = id\_username}
\end{center}

\subsection{Unified Class Architecture}
All six TC suites live in a single file \texttt{level2/TC-ALL.py}. A shared base class \texttt{\_BaseLevel2} owns the browser lifecycle, login, cookie-banner dismissal, and a session-recovery hook (\texttt{setUp()} probes \texttt{driver.current\_url} and re-creates the driver if it has died):

\begin{center}
\renewcommand{\arraystretch}{1.3}
\begin{tabular}{|l|l|c|}
\hline
\textbf{Suite class} & \textbf{Feature} & \textbf{Test count} \\ \hline
\texttt{TestCreateUserLevel2} & TC-001 -- Add User & 43 \\ \hline
\texttt{TestCreateCourseLevel2} & TC-002 -- Create Course & 27 \\ \hline
\texttt{TestAssignLevel2} & TC-003 -- Create Assignment & 27 \\ \hline
\texttt{TestGradeLevel2} & TC-004 -- Grade Assignment & 17 \\ \hline
\texttt{TestCalendarEventLevel2} & TC-005 -- Calendar Event & 28 \\ \hline
\texttt{TestQuizSetupLevel2} & TC-006 -- Quiz Setup & 27 \\ \hline
\multicolumn{2}{|r|}{\textbf{Total}} & \textbf{169} \\ \hline
\end{tabular}
\end{center}

\subsection{CSV Schema Example (Excerpt -- TC-001 Level~2)}
\begin{lstlisting}[language=, basicstyle=\ttfamily\scriptsize, frame=single, breaklines=true]
test_case_id,site_url,login_url_suffix,admin_username,admin_password,new_user_url,username,password,
firstname,lastname,email,expected_result,
username_locator_type,username_locator_value,
firstname_locator_type,firstname_locator_value,
lastname_locator_type,lastname_locator_value,
email_locator_type,email_locator_value,
save_btn_locator_type,save_btn_locator_value

TC-001-001,https://xuansang1234.moodlecloud.com,/login/index.php,
sang.truong2005@hcmut.edu.vn,Abcdxyz12@,https://xuansang1234.moodlecloud.com/user/editadvanced.php?id=-1,
usr001_uuuuuuuuuu,XuanSang12@,ffffffffff,llllllllll,usr001_nom@test.com,
success,id,id_username,id,id_firstname,id,id_lastname,id,id_email,
id,id_submitbutton
\end{lstlisting}

\subsection{Test Execution Summary -- Level 2}
A complete Level~2 run on the pinned environment yields:
\begin{center}
\fbox{\parbox{0.8\textwidth}{\centering
\texttt{====== 169 passed in 2323.73s (0:38:43) ======}
}}
\end{center}
i.e.\ all 169 parametric tests pass green; total wall-clock time $\approx$~39~minutes.

% ============================================================
% 5. NON-FUNCTIONAL TESTING
% ============================================================
\section{Non-functional Testing}

The project description requires \emph{at least one} non-functional requirement (NFR) per student (lecturer-adjusted from the original two), with each NFR described in terms of \textbf{testing type, testing approach, and testing tool}. With five team members and six functional features, the team implemented \textbf{six distinct NFR techniques} -- one per feature file -- covering a broad spectrum of quality attributes:
\begin{itemize}
    \item \textbf{Performance} -- load testing with \textbf{Locust} and SLA assertions (\texttt{TC-001.py}).
    \item \textbf{Security} -- passive HTTP probes with the \textbf{requests} library (\texttt{TC-002.py}).
    \item \textbf{Accessibility} -- WCAG 2.1 audit with \textbf{axe-selenium-python} (\texttt{TC-003.py}).
    \item \textbf{Reliability} -- repeated-request consistency checks (\texttt{TC-004.py}).
    \item \textbf{Compatibility} -- responsive-viewport rendering verification (\texttt{TC-005.py}).
    \item \textbf{Usability} -- keyboard-navigation and focus-indicator checks (\texttt{TC-006.py}).
\end{itemize}

Each NFR file targets one functional feature and runs as a standalone \texttt{python -m pytest TC-XXX.py} command. All six files reside in the \texttt{non\_functional/} directory. The project PDF's \emph{type / approach / tool} triad is documented for each technique below.

\subsection{NFR Catalogue}

\subsubsection{NFR-1 -- Performance: Load Testing (\texttt{TC-001.py})}
\textbf{Type:} Performance / concurrent load (throughput, p95 latency, failure rate).\\
\textbf{Approach:} Define a virtual user as a Python class deriving from \texttt{HttpUser}; the \texttt{@task} method loads the Add-User form under an authenticated admin session. Locust spawns virtual users in parallel and aggregates timing statistics in real time through its web UI on port~8089. An additional \texttt{unittest.TestCase} class measures login-page and Add-User form load times against SLA thresholds using \texttt{requests}~+~\texttt{time.time()}. To prevent Locust's gevent monkey-patch from racing pytest's selenium imports, the \texttt{HttpUser} subclass is only instantiated when the file is invoked outside pytest.\\
\textbf{Tool:} Locust 2.43 (\url{https://locust.io}).\\
\textbf{SLAs:} login page $\leq 5$~s; Add-User form load $\leq 8$~s.

\subsubsection{NFR-2 -- Security: Passive HTTP Probes (\texttt{TC-002.py})}
\textbf{Type:} Security / passive black-box probing of authentication, CSRF protection, XSS reflection, and PHP error disclosure.\\
\textbf{Approach:} A \texttt{requests.Session} is used to (1) confirm that anonymous access to the course-creation endpoint is blocked (redirects to login); (2) verify the login form embeds a CSRF \texttt{logintoken} and the authenticated session exposes a \texttt{sesskey} CSRF token; (3) POST an \texttt{<script>}-based XSS payload as the username and assert the raw script is never echoed back unescaped; (4) GET endpoints with deliberately malformed query parameters (\texttt{' OR 1=1--}, oversized strings, raw \texttt{<script>} in IDs) and assert the response body contains no PHP error signatures (\texttt{Fatal error}, \texttt{Call to undefined function}, \texttt{Stack trace:}, etc.). Daemon-free; runs against the live MoodleCloud instance without any prior setup.\\
\textbf{Tool:} \texttt{requests} 2.31 (Python standard ecosystem, no extra service).

\subsubsection{NFR-3 -- Accessibility: WCAG Audit with axe-core (\texttt{TC-003.py})}
\textbf{Type:} Accessibility / WCAG 2.1 conformance.\\
\textbf{Approach:} While a Selenium session has the assignment-creation form open, \texttt{axe-core} (the industry-standard accessibility engine) is injected into the live DOM through \texttt{axe-selenium-python}. The engine audits the page against $\sim$90~rules covering colour contrast, ARIA labelling, keyboard focus order, and heading hierarchy. The test fails if any violation of impact \emph{``critical''} or \emph{``serious''} is reported. A full JSON evidence report is persisted alongside the test.\\
\textbf{Tool:} \texttt{axe-selenium-python} 2.1 (wraps \texttt{axe-core} 4.x).

\subsubsection{NFR-4 -- Reliability: Repeated-Request Consistency (\texttt{TC-004.py})}
\textbf{Type:} Reliability / session stability and response consistency.\\
\textbf{Approach:} The grader page is loaded $N=5$ consecutive times within the same authenticated \texttt{requests.Session}. Each load asserts HTTP~200, the presence of key page markers, and session persistence (no redirect to login). Additionally, no single load may exceed $3\times$ the fastest observed response time (\texttt{DEGRADATION\_FACTOR~=~3.0}), ensuring the server does not silently degrade under repeated identical requests.\\
\textbf{Tool:} \texttt{requests} 2.31.

\subsubsection{NFR-5 -- Compatibility: Responsive Viewport Rendering (\texttt{TC-005.py})}
\textbf{Type:} Compatibility / responsive-design cross-device rendering.\\
\textbf{Approach:} A Selenium session opens the calendar month view and the browser window is resized to three standard breakpoints: Desktop (1920$\times$1080), Tablet (768$\times$1024), and Mobile (375$\times$812). At each viewport the test asserts that the calendar grid selector (\texttt{.calendarwrapper}, \texttt{table.calendartable}, or \texttt{.calendar-monthly-cell}) is visible and that no element overflows the viewport boundary, confirming the UI degrades gracefully at smaller widths.\\
\textbf{Tool:} Selenium 4 (window-resize API).

\subsubsection{NFR-6 -- Usability: Keyboard Navigation and Focus (\texttt{TC-006.py})}
\textbf{Type:} Usability / keyboard operability and error visibility.\\
\textbf{Approach:} A Selenium session opens the quiz-creation form and cycles through the form elements using the \texttt{Tab} key (up to \texttt{TAB\_LIMIT~=~80} presses) to verify that all interactive controls are keyboard-reachable, each focused element has a visible focus indicator (non-zero \texttt{outline} or \texttt{border} in computed style), and every control has an associated \texttt{<label>} element. A separate test submits an empty form and asserts that validation error messages appear in the DOM.\\
\textbf{Tool:} Selenium 4 (keyboard action API).

\subsection{NFR Coverage Matrix}

\begin{table}[H]
\centering
\renewcommand{\arraystretch}{1.3}
\begin{tabularx}{\textwidth}{|l|X|l|}
\hline
\textbf{File} & \textbf{Feature target} & \textbf{NFR technique} \\ \hline
\texttt{TC-001.py} & Admin Adds User              & Performance (Locust) \\ \hline
\texttt{TC-002.py} & Admin Creates Course         & Security (\texttt{requests}) \\ \hline
\texttt{TC-003.py} & Teacher Creates Assignment   & Accessibility (axe-core) \\ \hline
\texttt{TC-004.py} & Teacher Grades Student       & Reliability (\texttt{requests}) \\ \hline
\texttt{TC-005.py} & Admin Creates Calendar Event & Compatibility (Selenium) \\ \hline
\texttt{TC-006.py} & Teacher Sets Up a Quiz       & Usability (Selenium) \\ \hline
\end{tabularx}
\caption{NFR coverage -- one technique per functional feature}
\end{table}

% ============================================================
% 6. SUPPORTING TOOLING
% ============================================================
\section{Supporting Tooling}

\begin{tcolorbox}[colback=red!6!white, colframe=red!70!black, title=\textbf{Important -- Run Tests Manually}, fonttitle=\small\bfseries, left=4pt, right=4pt, top=4pt, bottom=4pt]
\small
All pytest / Selenium test commands must be executed \textbf{manually in your own terminal} (PowerShell, bash, or cmd).
Do \textbf{not} run them through an AI agent terminal (e.g.\ GitHub Copilot, Cursor, or any IDE agent runner).
Agent-managed shells may interfere with Chrome's DevTools Protocol, steal focus from the browser window, or terminate the WebDriver process mid-test -- causing spurious failures or an unrecoverable browser session.
\end{tcolorbox}

\subsection{Setup}
Install all dependencies with a single command (requires Python~3.9+ and Google Chrome; \texttt{webdriver-manager} auto-downloads the matching ChromeDriver):

\begin{lstlisting}[language=bash, basicstyle=\ttfamily\footnotesize, frame=single]
# macOS / Linux
python3 -m pip install -r requirements.txt

# Windows
pip install -r requirements.txt
\end{lstlisting}

\subsection{Per-Suite Run Commands}

\paragraph{Level 1 -- individual TCs}
\begin{lstlisting}[language=bash, basicstyle=\ttfamily\footnotesize, frame=single]
cd level1

python -m pytest TC-001_code.py -v   # TC-001  Admin Adds a New User         (43 cases)
python -m pytest TC-002_code.py -v   # TC-002  Admin Creates a New Course     (27 cases)
python -m pytest TC-003_code.py -v   # TC-003  Teacher Creates an Assignment  (27 cases)
python -m pytest TC-004_code.py -v   # TC-004  Teacher Grades an Assignment   (17 cases)
python -m pytest TC-005_code.py -v   # TC-005  User Creates a Calendar Event  (28 cases)
python -m pytest TC-006_code.py -v   # TC-006  Teacher Sets Up a Quiz         (27 cases)

python -m pytest . -v                # All Level 1 at once

# Run a single test case by ID
python -m pytest TC-001_code.py -v -k "TC_001_001"
\end{lstlisting}

\paragraph{Level 2 -- single file, one class per TC}
\begin{lstlisting}[language=bash, basicstyle=\ttfamily\footnotesize, frame=single]
cd level2

python -m pytest TC-ALL.py -v                                     # All 169 cases

python -m pytest TC-ALL.py -v -k "TestCreateUserLevel2"           # TC-001
python -m pytest TC-ALL.py -v -k "TestCreateCourseLevel2"         # TC-002
python -m pytest TC-ALL.py -v -k "TestAssignLevel2"               # TC-003
python -m pytest TC-ALL.py -v -k "TestGradeLevel2"                # TC-004
python -m pytest TC-ALL.py -v -k "TestCalendarEventLevel2"        # TC-005
python -m pytest TC-ALL.py -v -k "TestQuizSetupLevel2"            # TC-006

# Run a single test case by ID
python -m pytest TC-ALL.py -v -k "TC_004_002"
\end{lstlisting}

\paragraph{Non-functional tests}
\begin{lstlisting}[language=bash, basicstyle=\ttfamily\footnotesize, frame=single]
cd non_functional

python -m pytest TC-001.py -v   # Performance   -- SLA checks
python -m pytest TC-002.py -v   # Security      -- auth + CSRF + XSS probes
python -m pytest TC-003.py -v   # Accessibility -- axe WCAG audit
python -m pytest TC-004.py -v   # Reliability   -- repeated load consistency
python -m pytest TC-005.py -v   # Compatibility -- Desktop / Tablet / Mobile
python -m pytest TC-006.py -v   # Usability     -- keyboard nav + focus

python -m pytest . -v           # All NFR at once

# Locust load test (TC-001 Performance -- opens UI at http://localhost:8089)
locust -f TC-001.py
\end{lstlisting}

\paragraph{Run everything at once}
\begin{lstlisting}[language=bash, basicstyle=\ttfamily\footnotesize, frame=single]
# macOS / Linux
python3 -m pytest level1/ level2/ non_functional/ -v

# Windows
python -m pytest level1/ level2/ non_functional/ -v
\end{lstlisting}

Note: Locust load tests are interactive and are not included in the bulk pytest run above.
Use \texttt{run\_all.ps1 -Mode locust} or the \texttt{locust -f} command above to run them separately.

\subsection{Master Test Runners -- \texttt{run\_all.ps1} / \texttt{run\_all.sh}}
A single entry point orchestrates the entire test pipeline (install $\to$ sequenced suite execution). \texttt{run\_all.ps1} targets Windows PowerShell; \texttt{run\_all.sh} provides the equivalent for macOS / Linux bash. Both expose the same ten non-interactive modes plus an interactive menu:

\begin{lstlisting}[language=PowerShell, basicstyle=\ttfamily\footnotesize]
.\run_all.ps1                    # Interactive menu
.\run_all.ps1 -Mode setup        # Install dependencies, probe Chrome
.\run_all.ps1 -Mode smoke        # 1-test sanity check (~30 s)
.\run_all.ps1 -Mode level1       # All 169 Level 1 tests (~30 min)
.\run_all.ps1 -Mode level2       # All 169 Level 2 tests (~30 min)
.\run_all.ps1 -Mode nfr          # 6 NFR files (Performance/Security/Accessibility/Reliability/Compatibility/Usability)
.\run_all.ps1 -Mode locust       # Pick a Locust load-test file
.\run_all.ps1 -Mode tc -Tc 003   # Run TC-003 across level1+level2+NFR
.\run_all.ps1 -Mode cleanup      # Delete Moodle test data
.\run_all.ps1 -Mode all          # Everything end-to-end (~1 h 30 m)
\end{lstlisting}

The suite runs from a single \texttt{pip install -r requirements.txt} on a clean machine.

\subsection{Test-data Cleanup -- \texttt{cleanup\_moodle.py}}
Because the data-driven tests intentionally do \emph{not} clean up after themselves (sequential dependencies between rows -- e.g.\ \texttt{TC-001-034} validates that creating a user with username \texttt{usr001\_uuuuuuuuuu} fails as a duplicate, so TC-001-001 must have run first), accumulated test rows would eventually pollute the Moodle instance. The cleanup utility addresses this between full runs:

\begin{itemize}
    \item \textbf{Users} -- bulk delete on \texttt{/admin/user/user\_bulk.php} matching pattern \texttt{usr\textbackslash{}d+\_*}, mirroring the Katalon TC-001-Cleanup flow.
    \item \textbf{Courses} -- search-page enumeration matching pattern \texttt{fn*}~/~\texttt{sn*}, deleted through the two-step \texttt{/course/delete.php} confirmation flow.
    \item \textbf{Activities} -- per-course enumeration (course~10 for TC-003 assignments, course~12 for TC-006 quizzes) with names matching test patterns, deleted through a JavaScript \texttt{fetch}~+~\texttt{DOMParser} fast-path that submits the Moodle confirm form directly.
    \item \textbf{Calendar events} -- enumeration from \texttt{/calendar/view.php?view=upcoming}, deleted through \texttt{/calendar/delete.php}.
\end{itemize}

Protected entities (built-in \texttt{admin} user, the team's administrator account, courses 10 and 12 themselves) are excluded by a hard-coded allowlist.

\subsection{Reproducibility Manifest -- \texttt{requirements.txt}}
\begin{lstlisting}[language=, basicstyle=\ttfamily\footnotesize, frame=single]
selenium>=4.10           # WebDriver core
webdriver-manager>=4.0   # Auto-managed ChromeDriver
pandas>=2.0              # CSV / Excel loading
openpyxl>=3.1            # XLSX support
pytest>=7.4              # Test runner
locust>=2.20             # Performance load testing
requests>=2.31           # Passive HTTP security probes (NFR-2)
axe-selenium-python>=2.1.6     # Accessibility auditing
\end{lstlisting}
A single \texttt{pip install -r requirements.txt} restores the entire stack on a clean machine.

% ============================================================
% 7. TEST EXECUTION RESULTS
% ============================================================
\section{Test Execution Results}

\subsection{Aggregate Pass / Fail Summary}
The final clean run on the pinned environment produces the following results:

\begin{table}[H]
\centering
\renewcommand{\arraystretch}{1.3}
\begin{tabularx}{\textwidth}{|l|c|c|c|X|}
\hline
\textbf{Suite} & \textbf{Total} & \textbf{Pass} & \textbf{Fail} & \textbf{Wall-clock} \\ \hline
Level 1 -- TC-001 & 43 & 43 & 0 & $\sim$8 min \\ \hline
Level 1 -- TC-002 & 27 & 27 & 0 & $\sim$5 min \\ \hline
Level 1 -- TC-003 & 27 & 27 & 0 & $\sim$5 min \\ \hline
Level 1 -- TC-004 & 17 & 17 & 0 & $\sim$3 min \\ \hline
Level 1 -- TC-005 & 28 & 28 & 0 & $\sim$5 min \\ \hline
Level 1 -- TC-006 & 27 & 27 & 0 & $\sim$5 min \\ \hline
\textbf{Level 1 total} & \textbf{169} & \textbf{169} & \textbf{0} & $\sim$31 min \\ \hline
\textbf{Level 2 total} & \textbf{169} & \textbf{169} & \textbf{0} & 38 min 43 s \\ \hline
NFR-1 Performance (Locust/SLA) & -- & all pass & 0 & $\sim$1 min \\ \hline
NFR-2 Security (\texttt{requests}) & -- & all pass & 0 & $\sim$1 min \\ \hline
NFR-3 Accessibility (axe-core) & -- & all pass & 0 & $\sim$2 min \\ \hline
NFR-4 Reliability (\texttt{requests}) & -- & all pass & 0 & $\sim$1 min \\ \hline
NFR-5 Compatibility (Selenium) & -- & all pass & 0 & $\sim$3 min \\ \hline
NFR-6 Usability (Selenium) & -- & all pass & 0 & $\sim$3 min \\ \hline
\end{tabularx}
\caption{Final clean-run summary}
\end{table}

% ============================================================
% 8. CONCLUSION
% ============================================================

\section{Conclusion}
Project~\#3 successfully transformed the manual black-box test designs of Project~\#2 into a robust, reproducible, two-level data-driven automation suite covering all six target features (169 test methods per level, 338 in total), plus six non-functional test files each implementing a distinct quality-attribute technique: Performance (Locust), Security (\texttt{requests}), Accessibility (axe-core), Reliability, Compatibility, and Usability.

Throughout this project we gained practical experience in:
\begin{itemize}
    \item \textbf{Test architecture} -- balancing the trade-off between locator hard-coding (Level 1, simpler) and full CSV externalisation (Level 2, more flexible but more verbose CSV schemas).
    \item \textbf{Cross-browser reliability} -- using \texttt{webdriver-manager} to insulate the suite from ChromeDriver / Chrome version drift, and JavaScript-based DOM manipulation to bypass overlay interception, React-controlled inputs, and read-only attributes.
    \item \textbf{Non-functional methodology} -- integrating six distinct quality-attribute techniques (Performance, Security, Accessibility, Reliability, Compatibility, Usability) into a unified Python testing harness, and documenting each according to the project rubric's \emph{type / approach / tool} triad. Lightweight daemon-free approaches (\texttt{requests}, Selenium) were preferred throughout to keep the test suite runnable from a single \texttt{pip install} without external service dependencies.
    \item \textbf{Test-data lifecycle} -- recognising that sequential dependencies between data rows are intentional, and that cleanup must therefore run \emph{between} suite runs rather than \emph{after every test}.
    \item \textbf{Operational tooling} -- packaging a complex multi-stage test pipeline behind a single \texttt{run\_all.ps1} entry point, with a dedicated \texttt{-Mode tc} mode that runs one TC across all three directories at once and a \texttt{pytest.ini}-driven \texttt{TC-*} discovery convention that keeps the directory listing self-documenting.
\end{itemize}

The final clean run achieves \textbf{169/169 green on Level~1}, \textbf{169/169 green on Level~2}, and \textbf{all non-functional assertions passing} on the pinned environment, demonstrating that the Moodle LMS demo instance meets the documented functional, performance, security, accessibility, reliability, compatibility, and usability expectations of this academic test scope.

\end{document}
